% MA 214 Discussion 5 -- Module 1 synthesis and Quiz 1 review.
%
% Student handout. Page 1 is the Module 1 concept review; the question set runs from page 2.
% One data set, every Module 1 objective. Each question is tagged with the objective it tests,
% so students can see which one to go back to.
%
% Compile from inside this folder:  pdflatex Discussion5.tex

\documentclass[11pt]{article}
\usepackage[margin=1in]{geometry}
\usepackage{parskip}
\usepackage{fancyhdr}
\usepackage{array}
\usepackage{booktabs}
\usepackage{enumitem}
\usepackage{amsmath}
\usepackage{tabularx}
\usepackage{listings}

\pagestyle{fancy}
\fancyhf{}
\cfoot{\thepage}
\rhead{Discussion 5}
\lhead{MA 214: Applied Statistics}
\renewcommand{\headrulewidth}{.4mm}

\newcommand{\blankline}[1]{\underline{\hspace{#1}}}
\newcolumntype{L}{>{\raggedright\arraybackslash}X}

\lstset{
  basicstyle=\ttfamily\small,
  columns=fullflexible,
  frame=single,
  framerule=0.4pt,
  breaklines=true,
  showstringspaces=false,
  aboveskip=8pt,
  belowskip=8pt
}

% Section header: bold title over a hairline rule.
\newcommand{\parttitle}[1]{%
  \par\medskip\noindent
  {\large\bfseries #1}\par
  \vspace{-.5em}\noindent\rule{\linewidth}{0.4pt}\par\smallskip}

% Writing space.
\newcommand{\answerspace}[1]{\vspace{#1}}

% Objective tag printed after a question title.
\newcommand{\lo}[1]{{\normalfont\small (#1)}}

% Flush-left question list: the label runs in at the margin and the body wraps
% back to the margin, so nothing is pushed far to the right.
\newlist{questions}{enumerate}{1}
\setlist[questions]{label=\textbf{Question \arabic*.}, wide=0pt, itemsep=1.4em, topsep=.8em}

\newlist{parts}{enumerate}{1}
\setlist[parts]{label=\textbf{(\Alph*)}, leftmargin=1.6em, labelsep=.45em,
  itemsep=.7em, topsep=.5em}

\begin{document}

\noindent\textbf{Name:}\blankline{2.3in}\hfill\textbf{BUID:}\blankline{1.6in}

\begin{center}
  {\Large\bfseries Discussion 5: All of Module 1 on One Data Set}
\end{center}

\noindent\textbf{Goals:} run every Module 1 tool on a single data set, and find out which
objective you still need to go back to before the quiz. Each question is labelled with the
objective it tests.

\parttitle{Part 1: Module 1 Concept Review}

{\small

\noindent The whole module on one page. Everything in Part 2 can be answered from it.

\begin{enumerate}[label=\textbf{\arabic*.}, wide=0pt, itemsep=.45em, topsep=.5em]

\item \textbf{M1.L01 Models, likelihood, and maximum likelihood.} A statistical model gives the
data a \emph{distribution} with unknown parameters. To estimate: write the likelihood $L(\theta)$,
take logs, set $\ell'(\theta)=0$, solve. For $y$ successes in $n$ trials,
$\hat\theta = y/n$; for exponential waits, $\hat\lambda = 1/\bar{x}$. For the linear model,
maximizing the likelihood under \emph{normal errors with constant variance} is exactly
minimizing $\sum e_i^2$, which is why least squares is a maximum likelihood estimator.

\item \textbf{M1.L02 The linear model and its coefficients.} The model is
$Y_i \mid x_i \sim \text{Normal}(\beta_0 + \beta_1 x_{1i} + \cdots, \sigma^2)$, independently, so
it assumes linearity, independence, constant variance, and normality. In multiple regression each
coefficient is the predicted change in $y$ per one-unit increase in that predictor
\textbf{holding the others fixed}. Always report units. $R^2$ is the fraction of variability in
$y$ explained by the model; it says nothing about whether the model is \emph{right}, and nothing
about steepness.

\item \textbf{M1.L03 Using a model: collinearity and extrapolation.} When predictors carry nearly
the same information, standard errors inflate and each can look useless while the pair matters.
Interpretation breaks; prediction survives. Predicting inside the observed range of the predictors
is \textbf{interpolation}; outside it is \textbf{extrapolation}, and the arithmetic never warns you.

\item \textbf{M1.L04 Model selection.} With $p$ predictors and $n$ observations,
\[
R^2_{\text{adj}} = 1 - (1-R^2)\,\frac{n-1}{n-p-1},
\qquad
\text{AIC} = -2\,\ell(\hat\theta) + 2k \;\; (\text{lower is better}).
\]
$R^2$ can only rise when a predictor is added, so it cannot compare models of different sizes.
Adjusted $R^2$ and AIC both charge for complexity and can move either way. AIC differences under
about 2 are not meaningful.

\item \textbf{M1.L05 Indicator variables.} A categorical predictor with $k$ levels enters as
$k-1$ indicators. The omitted level is the \textbf{reference}: the intercept describes it, and
each indicator coefficient is a \emph{difference} from it, holding the other predictors fixed.
Changing the reference changes the coefficients but not the fitted values.

\item \textbf{M1.L06 Logistic regression.} For a binary response, model the log odds as linear:
$\eta = \log\frac{p}{1-p} = \beta_0 + \beta_1x_1 + \cdots$, so
$p = e^{\eta}/(1+e^{\eta})$. Fit by maximum likelihood, not least squares. Exponentiating gives an
\textbf{odds ratio}: $e^{\beta}$ per one unit, $e^{c\beta}$ per $c$ units. To judge a classifier,
put accuracy next to the \textbf{base rate}, and report sensitivity and specificity too.

\end{enumerate}

}

\newpage

\parttitle{Part 2: Question Set}

\noindent\textbf{Scenario.} A city library system has $n = 60$ branches. For each branch you have
\texttt{hours} (weekly open hours, ranging 21 to 70), \texttt{staff} (staff FTE),
\texttt{type} (Branch, Central, or Kiosk), \texttt{visits} (annual visits in thousands, ranging
18.4 to 153.3), and \texttt{met} (whether the branch met its annual program target, $1$ = yes).
\textbf{18 of the 60 branches met the target.} In this data set \texttt{hours} and \texttt{staff}
are correlated at $0.77$.

\begin{questions}

%%%%%%%%%%%%%%%%%%%%%%%%%%%%%%%%%%%%%%
\item \textbf{Model and estimate.} \lo{M1.L01}

\begin{parts}
\item Treat the number of branches meeting the target as $Y \sim \text{Binomial}(60, \theta)$.
Write $L(\theta)$ and $\ell(\theta)$, then find $\hat\theta$ and report it as a number.
\answerspace{5em}

\item State the assumption about branches that the binomial model requires, and say whether you
believe it here.
\answerspace{3.5em}

\item You now fit a linear model for \texttt{visits} by least squares. Which model assumption
makes least squares the \emph{maximum likelihood} estimator, and what would still be true about
$S_{xy}/S_{xx}$ if that assumption failed?
\answerspace{4em}
\end{parts}

%%%%%%%%%%%%%%%%%%%%%%%%%%%%%%%%%%%%%%
\item \textbf{Reading the coefficient table.} \lo{M1.L02, M1.L05} Regressing \texttt{visits} on
\texttt{hours} and \texttt{type}:

\begin{center}
\renewcommand{\arraystretch}{1.2}
\begin{tabular}{lccc}
\toprule
 & \textbf{Estimate} & \textbf{Std. Error} & \textbf{$p$-value} \\
\midrule
(Intercept) & $1.791$ & 3.090 & 0.564 \\
\texttt{hours} & $1.506$ & 0.065 & $<0.001$ \\
\texttt{type}Central & $40.826$ & 2.212 & $<0.001$ \\
\texttt{type}Kiosk & $-11.802$ & 2.043 & $<0.001$ \\
\bottomrule
\end{tabular}
\qquad $R^2 = 0.951$
\end{center}

\begin{parts}
\item Which level is the reference, and how do you know from the table alone?
\answerspace{2.5em}

\item Interpret $1.506$ in context, with units and the conditional phrase.
\answerspace{3.5em}

\item Interpret $-11.802$. Your sentence must name the comparison group.
\answerspace{3.5em}

\item Predict annual visits for each type at \texttt{hours} $= 40$.

\begin{center}
\renewcommand{\arraystretch}{1.7}
\begin{tabular}{lcc}
\toprule
\textbf{Type} & \textbf{Arithmetic} & \textbf{Predicted \texttt{visits}} \\
\midrule
Branch & \blankline{2.1in} & \blankline{.9in} \\
Central & \blankline{2.1in} & \blankline{.9in} \\
Kiosk & \blankline{2.1in} & \blankline{.9in} \\
\bottomrule
\end{tabular}
\end{center}
\end{parts}

%%%%%%%%%%%%%%%%%%%%%%%%%%%%%%%%%%%%%%
\item \textbf{Choosing and misusing the model.} \lo{M1.L03, M1.L04} Three models on the same 60
branches.

\begin{center}
\renewcommand{\arraystretch}{1.3}
\begin{tabular}{llccccc}
\toprule
\textbf{Model} & \textbf{Predictors} & \textbf{$p$} & \textbf{$R^2$} & \textbf{$R^2_{\text{adj}}$}
 & \textbf{AIC} & \textbf{SE of \texttt{hours}} \\
\midrule
1 & \texttt{hours} & 1 & 0.4965 & 0.4879 & 541.06 & 0.204 \\
2 & $+$ \texttt{type} & 3 & 0.9512 & 0.9486 & 405.02 & 0.065 \\
3 & $+$ \texttt{staff} & 4 & 0.9537 & 0.9503 & 403.94 & 0.102 \\
\bottomrule
\end{tabular}
\end{center}

\noindent In Model 3 the \texttt{staff} coefficient is $1.406$ with $p = 0.094$.

\begin{parts}
\item Model 1 to Model 2 raises $R^2$ from 0.497 to 0.951. What does that say about leaving
\texttt{type} out?
\answerspace{3em}

\item Model 3 has the best $R^2_{\text{adj}}$ \emph{and} the best AIC, yet its \texttt{staff}
coefficient has $p = 0.094$ and the standard error of \texttt{hours} rose from 0.065 to 0.102.
Explain both facts, and say which model you would report and why.
\answerspace{5.5em}

\item Using Model 2, someone predicts visits for a branch open $120$ hours per week and gets
$182.5$ thousand. Give two separate reasons to distrust that number.
\answerspace{4.5em}

\item Would the collinearity between \texttt{hours} and \texttt{staff} damage Model 3's
\emph{predictions}? Answer in one sentence.
\answerspace{3em}
\end{parts}

%%%%%%%%%%%%%%%%%%%%%%%%%%%%%%%%%%%%%%
\item \textbf{A binary outcome.} \lo{M1.L06} Now model \texttt{met} on \texttt{visits} by logistic
regression:

\begin{center}
\renewcommand{\arraystretch}{1.2}
\begin{tabular}{lccc}
\toprule
 & \textbf{Estimate} & \textbf{Std. Error} & \textbf{$p$-value} \\
\midrule
(Intercept) & $-6.990$ & 1.727 & $<0.001$ \\
\texttt{visits} & $0.0745$ & 0.0197 & $<0.001$ \\
\bottomrule
\end{tabular}
\end{center}

\noindent You may use: $e^{0.0745} = 1.077$, \; $e^{0.745} = 2.107$, \;
$e^{-2.518} = 0.081$, \; $e^{0.464} = 1.590$, \; $e^{0.36} = 1.433$, \; $e^{1.13} = 3.096$.

\begin{parts}
\item Give the odds ratio for a one-unit (one thousand visits) increase, and for a ten-unit (ten
thousand visits) increase. Interpret the second one in a sentence a library director would
understand.
\answerspace{4.5em}

\item The 95\% interval for the \texttt{visits} coefficient is $[0.036,\; 0.113]$. Give the
interval for the ten-thousand-visit odds ratio, and say what it means that the interval excludes 1.
\answerspace{4em}

\item Complete the table. \texttt{visits} is in thousands.

\begin{center}
\renewcommand{\arraystretch}{1.6}
\begin{tabular}{cccc}
\toprule
\texttt{visits} & \textbf{Log odds $\eta$} & \textbf{Odds} & \textbf{Probability} \\
\midrule
60 & $-2.518$ & \blankline{1in} & \blankline{1in} \\
100 & $0.464$ & \blankline{1in} & \blankline{1in} \\
\bottomrule
\end{tabular}
\end{center}

\item At a threshold of $0.5$ this model classifies $78.3\%$ of branches correctly. Before calling
that good, what number must you put beside it, and what is that number here?
\answerspace{3.5em}
\end{parts}

%%%%%%%%%%%%%%%%%%%%%%%%%%%%%%%%%%%%%%
\item \textbf{Referee the report.} \lo{all of Module 1} A classmate submits the conclusions below.
Each numbered sentence is either defensible or not. Mark each, and rewrite the ones that are not.

\begin{enumerate}[label=\textbf{S\arabic*.}, leftmargin=2.2em, itemsep=.3em, topsep=.4em]
\item ``Our model explains 95.1\% of the variability in visits, so it is the correct model for
library visits.''
\item ``Each extra weekly open hour causes about 1.5 thousand more annual visits.''
\item ``Central branches receive 40.8 thousand visits per year.''
\item ``Because \texttt{staff} has $p = 0.094$, staffing does not affect visits.''
\item ``A branch open 120 hours per week would receive about 183 thousand visits.''
\item ``Our logistic model classifies 78\% of branches correctly, so it reliably predicts which
branches will meet their target.''
\end{enumerate}

\begin{center}
\renewcommand{\arraystretch}{2.1}
\begin{tabularx}{\linewidth}{ccL}
\toprule
\textbf{Sentence} & \textbf{Defensible?} & \textbf{If not, what is wrong and what should it say} \\
\midrule
S1 & & \\
S2 & & \\
S3 & & \\
S4 & & \\
S5 & & \\
S6 & & \\
\bottomrule
\end{tabularx}
\end{center}

\end{questions}

\end{document}
